\documentclass[]{article}

\usepackage{amsmath}
\usepackage{multirow}
\usepackage{natbib}
\usepackage{graphicx} 


\begin{document}

\subsection{Dataset and simulation setup}
The training and evaluation data were generated from N-body simulations of 50-particle virialized Plummer spheres. The Plummer model provides a simple, spherically symmetric density profile commonly used to represent star clusters. A total of 100 independent simulations were run using a standard leapfrog integrator with a fixed timestep of $dt=0.01$ for a total duration of $T=5.0$ in simulation units. A gravitational softening length of $\epsilon=0.01$ was used to prevent numerical divergences during close particle encounters. From these simulations, 80 were allocated for training and 20 were held out for testing. To evaluate the model's generalization capabilities, additional test sets were generated for systems with $N=25$ and $N=100$ particles, following the same initialization and integration procedure.

To augment the training data and enforce physical symmetries, each batch of particle coordinates and velocities was subjected to a random 3D rotation and translation before being fed into the network. This ensures the model learns the inherent Galilean invariance of the gravitational N-body problem.

\subsection{Hamiltonian graph neural network architecture}
Our approach is centered on learning the system's scalar Hamiltonian, $H(q, p)$, rather than directly regressing the vector forces. The Hamiltonian is separated into kinetic and potential energy components, $H(q, p) = T(p) + U(q)$. The kinetic energy is defined analytically as $T(p) = \sum_{i=1}^{N} \frac{p_i^2}{2m_i}$, where we assume unit masses ($m_i=1$) for all particles.

The potential energy, $U(q)$, is parameterized by a permutation-invariant graph neural network (GNN). The GNN models the potential as a sum over pairwise interactions, an architecture well-suited for particle systems. The total potential energy is given by:
\begin{equation}
    U(q) = \frac{1}{N} \sum_{i<j} \phi(d_{ij})
\end{equation}
where $\phi$ is a shared multi-layer perceptron (MLP) with SiLU activation functions that maps the pairwise distance between particles $i$ and $j$ to a scalar potential. The $1/N$ scaling factor ensures the model's output remains stable across different particle counts.

A key physical prior is incorporated directly into the network architecture by using the softened pairwise distance as the input to the MLP:
\begin{equation}
    d_{ij} = \sqrt{\|q_i - q_j\|^2 + \epsilon^2}
\end{equation}
where $\epsilon=0.01$ is the same gravitational softening length used in the ground-truth data generation. This provides a strong inductive bias and stabilizes the learning process by avoiding the singularity of the pure Newtonian potential. The forces are then derived from the learned potential using automatic differentiation, which guarantees that the resulting force field is conservative (curl-free):
\begin{equation}
    \dot{p}_i = -\frac{\partial H}{\partial q_i} = -\frac{\partial U}{\partial q_i}
\end{equation}

\subsection{Symplectic integration and training}
To ensure the learned dynamics preserve the geometric structure of the phase space, we embed a differentiable second-order symplectic integrator directly into the training loop. We use the "kick-drift-kick" formulation of the leapfrog algorithm. A single integration step from time $t_n$ to $t_{n+1}$ with timestep $\Delta t$ proceeds as follows:
\begin{enumerate}
    \item \textbf{Kick:} Update momenta by a half-step: $p_{n+1/2} = p_n - \frac{\Delta t}{2} \nabla_q U(q_n)$.
    \item \textbf{Drift:} Update positions by a full step: $q_{n+1} = q_n + \Delta t \frac{p_{n+1/2}}{m}$.
    \item \textbf{Kick:} Update momenta by a final half-step: $p_{n+1} = p_{n+1/2} - \frac{\Delta t}{2} \nabla_q U(q_{n+1})$.
\end{enumerate}
The entire integration sequence is implemented using differentiable operations, allowing gradients to be backpropagated through time. The model is trained to predict the state at a future time by unrolling this integrator over multiple steps. Specifically, we use a sliding window approach, minimizing the Mean Squared Error (MSE) between the predicted state at $t_n + 50 \cdot dt$ and the ground-truth state, starting from the state at $t_n$.

To manage the multi-scale nature of the Plummer sphere, we employed a two-stage curriculum learning strategy. In the first stage, the loss was masked to include only particles outside the dense core (radial position $r > 0.5b$, where $b=1$ is the scale radius), allowing the model to first learn the smoother, long-range potential. In the second stage, the mask was removed to allow the model to refine its handling of high-density, short-range interactions within the core. To further stabilize training and prevent energy drift, a regularization term was added to the loss function: $\lambda \mathrm{E}[|H_{\mathrm{pred}} - H_{\mathrm{initial}}|^2]$, with $\lambda = 0.001$. The network was trained using the Adam optimizer.

\subsection{Evaluation metrics}
The performance of the trained emulator was assessed on several criteria, focusing on both trajectory accuracy and the preservation of fundamental physical symmetries.
\begin{itemize}
    \item \textbf{Trajectory Reconstruction:} The primary accuracy metric was the Mean Squared Error (MSE) in both position and velocity between the emulated trajectories and the ground-truth simulations over the full time interval $T=5.0$. This was evaluated on the held-out test sets for $N=25, 50,$ and $100$ particles to assess zero-shot generalization.
    \item \textbf{Hamiltonian Conservation:} To verify long-term energy stability, we measured the Hamiltonian deviation, $\Delta H(t) = |H(q(t), p(t)) - H(q(0), p(0))|$, throughout the integration. A symplectic model should exhibit bounded, oscillatory energy error, in contrast to the secular energy drift seen in non-symplectic methods.
    \item \textbf{Time-Reversibility:} We tested the time-reversal symmetry of the learned dynamics by integrating a system forward from its initial state at $t=0$ to the final state at $t=5.0$, and then integrating it backward in time (using a negative timestep) to $t=0$. The reversibility error was quantified as the final Euclidean distance between the backward-integrated state and the original initial state.
    \item \textbf{Phase-Space Volume Preservation:} According to Liouville's theorem, a Hamiltonian flow is volume-preserving in phase space. This implies that the Jacobian of the one-step flow map, $M = \frac{\partial(q_{n+1}, p_{n+1})}{\partial(q_n, p_n)}$, must have a determinant of one. We numerically computed $\det(M)$ for our learned integrator to confirm that it acts as a canonical transformation, a rigorous test of its symplectic nature.
\end{itemize}

\end{document}
                